% Source: https://artofproblemsolving.com/community/c3h1930276_help_needed_in_problems_of_varied_topics
% Stable identity: aops-topic-1930276-q17
Let $(x,y)$ be a solution. Since $\dfrac{1}{x}=\dfrac{1}{n}+\dfrac{1}{y}>\dfrac{1}{n}$, we have $x<n$; write $x=n-a$ with $1\le a\le n-1$. Then
\[
\dfrac{1}{y}=\dfrac{1}{n-a}-\dfrac{1}{n}=\dfrac{a}{n(n-a)},\qquad y=\dfrac{n(n-a)}{a}=\dfrac{n^2}{a}-n .
\]
So $y$ is a positive integer exactly when $a\mid n^2$ (and then $y>0$ because $a<n$). Hence the solutions correspond one-to-one to the divisors $a$ of $n^2$ with $1\le a\le n-1$.
\begin{itemize}
\item $n=1$: there is no such $a$, so there are no solutions.
\item $n=p$ prime: the divisors of $p^2$ are $1,p,p^2$, and only $a=1$ is less than $p$. So there is exactly one solution, $(x,y)=(p-1,\,p(p-1))$.
\item $n$ composite: $n=rs$ with $1<r<n$, and both $a=1$ and $a=r$ are divisors of $n^2$ less than $n$, so there are at least two solutions.
\end{itemize}
Therefore there is exactly one solution if and only if $n$ is prime. $\blacksquare$
